\documentclass[11pt,a4paper]{exam}
\usepackage[T1]{fontenc}
\usepackage{longtable}
\usepackage{graphicx}

\usepackage{hyperref}
%\urlstyle{same}  % don't use monospace font for urls
\setlength{\parindent}{0pt}
\setlength{\parskip}{6pt plus 2pt minus 1pt}

\usepackage{geometry}
\usepackage{multicol}
\usepackage{relsize}
\usepackage{pgf}
\usepackage{fancyvrb}

\geometry{
  %includeheadfoot,
  margin=2.54cm
}
\title{Worksheet: Java Generics and Reflection (introspection)}
\author{}
\date{}

\fvset{tabsize=2}
\fvset{fontsize=\relsize{-1}}
\fvset{frame=single}

\begin{document}

\maketitle

This worksheet reinforces your existing knowledge of Java Generics and Reflection. These \emph{techniques} are
commonly used in the technologies we will be examining in the rest of the module. 

\subsection*{Required:}

You should ensure you are using (at least) version 7 of the JDK.

\subsection*{Preparation:}

Create the following \texttt{Storage}, \texttt{BankAccount} and \texttt{Driver} classes in a directory of your choice.

\VerbatimInput{src/Storage.java}
\VerbatimInput{src/BankAccount.java}
\VerbatimInput{src/Driver.java}

Add the following code snippet to your \texttt{Driver} class \texttt{main} method, 
creating two different storage
objects with two different type specialisations:

\begin{Verbatim}
  Storage<BankAccount>  aStorage = new Storage<>();
  Storage<String>       sStorage = new Storage<>();
\end{Verbatim}

\begin{questions}

\question What are the reasons for using generics here? 

\begin{solution}
Generics are used to introduce two different types into a function	
\end{solution}

\item What are the benefits? 

\begin{solution}
\begin{itemize}
\item They enable programmers to implement generic algorithms. 
\item Programmers can implement generic algorithms that work on collections of different types, that can be customised, and are type-safe and easier to read
\end{itemize}
\end{solution}

\item\label{ques}
Now add the following code to your \texttt{Driver} class:
\begin{Verbatim}
  Class baCls = BankAccount.class;
  try {
    Object myAccount =  baCls.newInstance();
    aStorage.setValue(myAccount);
    
    // Deposit
    myAccount.deposit(15);
  }
  catch ( InstantiationException e ) {
    // ...
  }
  catch ( IllegalAccessException e ) {
    // ...
  }
\end{Verbatim}

\noindent Compile and analyse the compiler output. \\
What is the cause of the problem reported by the compiler, if any?

\begin{solution}
The problem is caused class \verb!Storage! which can't be applied to the
given type \verb!aStorage.setValue(myAccount);!
\end{solution}

\item
Now replace:
\begin{Verbatim}
    Object myAccount =  baCls.newInstance();
\end{Verbatim}
with
\begin{Verbatim}
    BankAccount myAccount =  baCls.newInstance();
\end{Verbatim}

\noindent How does this affect the compilation process? \\
What is the problem, if any? \\
What does the \texttt{myAccount} variable hold when the code is executed?\\
Decide whether your diagnosis from question (~\ref{ques}) was correct.

\begin{solution}
\begin{itemize}
\item The class will still not compile.
\item The problem is related to incompatible types, namely:
\begin{Verbatim}
            BankAccount myAccount = baCls.newInstance();
\end{Verbatim}
\item The code can’t be executed at this stage because of the compile error 
\item A dynamic cast is required.
\end{itemize}	
\end{solution}

\item
Now add an explicit dynamic cast:
\begin{Verbatim}
    BankAccount myAccount =  (BankAccount) baCls.newInstance();
\end{Verbatim}

\noindent What does the dynamic cast do here? \\
Is it the compiler that performs the cast operation or the Java runtime environment (JVM)? \\
Is this code safe?

\begin{solution}
We need to convert the new object to the same type \verb!BankAccount! 
so that we are able to obtain the new balance on the account. \\
The cast operation is performed by JVM that is why is called a \emph{dynamic} cast.\\
It is not safe to perform the comparison using a static cast due to 
location where the cast is performed. 	
\end{solution}

\item
Now replace your initial declaration:
\begin{Verbatim}
  Class baCls = BankAccount.class;
\end{Verbatim}
with
\begin{Verbatim}
  Class<BankAccount> baCls = BankAccount.class;
\end{Verbatim}

\noindent
Explain the compiler output? \\
Are there errors? \\
What is the reason? \\
What does it say about the role of generics?

\begin{solution}
All the code should compile correctly without errors, this is because all the objects 
have been defined correctly with the aid of a dynamic cast.\\
The role of generics in Java is solely a compile-time effect due to \emph{type erasure}.
\end{solution}

\item
Now add:
\begin{Verbatim}
  System.out.println( aStorage.getValue().showBalance() );

  if( aStorage.getClass() == sStorage.getClass() ) {
    System.out.println( "EQUAL" );
  } else {
    System.out.println( "NOT EQUAL" );
  }
\end{Verbatim}
\noindent
What is the run-time output? \\
Explain why you get such output and how does this relate to generics and their use with reflective 
instantiation of objects?

\begin{solution}
\begin{Verbatim}
115.0
EQUAL
BUILD SUCCESSFUL (total time: 0 seconds)
\end{Verbatim}

We obtain the number \verb!115! because this is now the new balance on the account.  
The initial balance was \verb!100! held by the \verb!BankAccount! object. When this object 
is called by the main method of the \verb!Driver! class we then added \verb!15! 
to the current balance.

That is the mechanics of the code --- generics enables addition compile time constraints to be 
applied to the types but, of course, this does not apply to runtime.  
	
\end{solution}

\end{questions}

\end{document}